% appB —— 附录 B $\CM$ 上的柯西--黎曼方程（原文 wcf-paper.tex 3395–3528 行，PDF 56–57 页；共 7 段、17 条编号公式 (B.1)–(B.17)）
% 记号纪律：$\CM$ = 三维模空间（总空间），$\vm$ = 四维库仑分支（底空间），全块不混。
% 交叉引用：第 5 节 = §5 Mutually non-local corrections（ch05a.tex \label{sec:constructing-moduli}）；式 (3.2) = 一般复结构 $J^{\pz}$（ch03.tex \label{eq:Jzeta}）。合稿时改回 \ref/\eqref。
\section{$\CM$ 上的柯西--黎曼方程}
\label{app:holomorphy}

本附录说明：$(\CM,g)$ 上关于复结构 $J^\pz$ 的柯西--黎曼方程，如何可以改写成 $\vm$ 上某个联络的平坦性方程，而该联络对 $\zeta$ 的依赖非常简单。我们并不假定 $g$ 来自我们在第 5 节所描述的构造；相反，我们只用一些一般事实，它们来自把 $(\CM,g)$ 认定为该规范理论在 $\IR^3\times S^1$ 上的模空间。% TODO 合稿改 \ref{sec:constructing-moduli}

下面，对每个 $\zeta\in\IC^\times$，我们考虑 $\CM$ 上关于复结构 $J^\pz$ 的柯西--黎曼方程
\begin{equation} \label{eq:cauchy-riemann}
\bar\partial f = 0
\end{equation}
我们将把这些方程改写成如下形式
\begin{align}
\partial_{u^i} f &= \CA_{u^i} f, \label{eq:cr-connection-1} \\
\partial_{\bar{u}^{\bar i}} f &= \CA_{{\bar u}^{\bar i}} f, \label{eq:cr-connection-2}
\end{align}
其中 $\CA_{u^i}$ 与 $\CA_{{\bar u}^{\bar i}}$ 是沿环面纤维作用的一阶微分算子（于是在 $\CM$ 的坐标 $(u,\bar u,\theta)$ 下，它们只涉及对 $\theta$ 的导数）；此外，它们还以简单的方式依赖于 $\zeta$：
\begin{align}
\CA_{u^i} &= \frac{1}{\zeta} \CA^{(-1)}_{u^i} + \CA^{(0)}_{u^i}, \\
\CA_{\bar u^{\bar i}} &= \CA_{\bar u^{\bar i}}^{(0)} + \zeta \CA_{\bar u^{\bar i}}^{(1)},
\end{align}
其中诸 $\CA^{(-1)}_{u^i}$ 在每一点都线性无关，诸 $\CA^{(1)}_{\bar{u}^{\bar{i}}}$ 亦然。

我们先把 \eqref{eq:cauchy-riemann} 改写成
\begin{equation} \label{eq:holomorphy}
(1 - i J^{(\zeta,\bar\zeta)}) df = 0.
\end{equation}
若把 $\zeta$ 与 $\bar\zeta$ 看作彼此独立的复变量，那么这个方程其实与 $\bar\zeta$ 无关。为看出这一点，只需在一个固定的 $\zeta$ 处论证即可，例如取 $\zeta=0$。将一般的复结构 (3.2) 特殊化到 $\zeta=0$，% TODO 合稿改 \eqref{eq:Jzeta}
\begin{equation}
J^{(\zeta=0,\bar\zeta)} = J_3 + i \bar\zeta J_+
\end{equation}
其中我们引入了 $J_+ = J_1 + i J_2$。其次注意 $J_+ J_3 = i J_+$，故 $J_+ (J_3 - i) = 0$，于是 $J_+$ 湮灭 $J_3$ 的 $-i$ 特征子空间。这样我们就证明了：$J^{(\zeta=0,\bar\zeta)}$ 与 $J_3$ 共享一个 $n$ 维的、特征值为 $-i$ 的特征子空间。为完成论证，我们还想知道：$J^{(\zeta=0,\bar\zeta)}$ 没有\emph{其他}特征值为 $-i$ 的特征向量。为此我们作一个类似的论证：固定 $\bar\zeta$ 而让 $\zeta$ 变动；这样得到 $J^{(\zeta=0,\bar\zeta)}$ 的 $n$ 个特征值为 $+i$ 的特征向量。于是由维数计数，已没有多余的余地。最终我们看到：$J^{(\zeta,\bar\zeta)}$ 的 $-i$ 特征子空间如所愿地与 $\bar\zeta$ 无关。

于是，在研究柯西--黎曼方程 \eqref{eq:holomorphy} 时，我们可以任取一个方便的 $\bar\zeta$。既然我们想弄清 \eqref{eq:holomorphy} 用 $\vm$ 上的 SW 环面纤维化来表达时是什么样子，自然就取 $\bar\zeta=0$；把它代入 \eqref{eq:holomorphy}，得
\begin{equation} \label{eq:hol-specialized}
(1 - i J_3 - \zeta J_-) df = 0.
\end{equation}
我们假定已给定复辛流形 $(\CM, J_3, \omega_+)$ 与 SW 环面纤维化 $(\CM, J_3^\sf, \omega_+^\sf = - \frac{1}{4 \pi} da^I \wedge dz_I)$ 之间的一个等同。（\cite{Seiberg:1996nz} 中论证过：借助规范理论与引力的弱耦合，这样的等同至少对 $J_3$ 应当存在；类似的论证表明，对 $\omega_+$ 也存在这样的等同。）然后，用与环面纤维相切的向量场 $\frac{\partial}{\partial \bar{z}_I}$ 去缩并 \eqref{eq:hol-specialized}，得到
\begin{equation}
2 \partial_{\bar{z}_I} f - \zeta \left( \partial_{\bar{z}_I} \cdot J_- df \right) = 0.
\end{equation}
为处理第二项，我们利用
\begin{equation}
J_- = g^{-1} \omega_- =  \frac{1}{4 \pi} g^{-1}(d\bar{a}^{\bar I}
\wedge d\bar{z}_{\bar I})
\end{equation}
再乘以 $4 \pi / \zeta$，得到
\begin{equation}
g^{-1}(df, d \bar{a}^I) =  \frac{8 \pi}{\zeta} \partial_{\bar{z}_I} f.
\end{equation}
这已经差不多是我们所要的形式 \eqref{eq:cr-connection-1}，但还不完全是：$g^{-1}(df, d \bar{a}^I)$ 是若干作用于 $f$ 的导数算子的混合。我们想作一次基变换，以便提取出关于 $\frac{\partial f}{\partial a^I}$ 的方程。为此，我们考虑 $g$ 在垂直于 $\CM_u$ 的水平子空间上的限制；把它写成 $g = h_{\bar{I} J} d\bar{a}^{\bar I} da^J$。再乘以 $h_{\bar{I} J}$，便得到
\begin{equation}
h_{\bar{I} J} g^{-1}(df, d \bar{a}^I) =  \frac{8 \pi}{\zeta}
h_{\bar{I} J} \partial_{\bar{z}_I} f.
\end{equation}
现在考虑 $f$ 只依赖于底空间坐标 $(a, \bar{a})$ 的特殊情形。在一点处取一组正交归一基，我们便看出 $g^{-1}(df, d \bar{a}^I) = (h^{-1})^{\bar{I} J} \partial_{a^J} f$。这意味着，对一般的 $f$，左端可以写成
\begin{equation}
\left( \partial_{a_J} - \CA^J \right) f
\end{equation}
其中 $\CA^J$ 是只在纤维方向上作用的微分算子。（更内蕴地说，完整的联络算子 $\partial_{a_J} - \CA^J$ 就是 $f$ 沿向量场 $\partial_{a_J}$ 从 $\vm$ 到 $\CM$ 的水平提升所作的导数。）

总之，我们已经得到
\begin{equation}
\frac{\partial}{\partial a_J} f = \frac{8 \pi}{\zeta} h_{\bar{I} J}
\frac{\partial f}{\partial \bar{z}_I} + \CA^J f,
\end{equation}
这正是所要的形式 \eqref{eq:cr-connection-1}。一个完全相同的论证（从 $\bar\zeta = \infty$ 而非 $\bar\zeta = 0$ 出发）给出共轭方程 \eqref{eq:cr-connection-2}。

我们在此发现的结构，与 \cite{Cecotti:1991me} 中引入的「改进的联络」非常相似。为最清楚地看出这种相似性，引入 $\vm$ 上的一个无穷维丛 $V$，使得 $V$ 在 $u \in \vm$ 处的纤维恰好是环面 $\CM_u$ 上实解析函数的空间，
\begin{equation}
V_u = C^\omega(\CM_u).
\end{equation}
于是，整个 $\CM$ 上的实解析复值函数 $f$，等价于 $V$ 在 $\vm$ 上的一个实解析截面。这样，上面所得到的结果就是：柯西--黎曼方程可以看作 $V$ 中一族具有特定形式的单参数联络的平坦性方程。这些联络的平坦性，是 $\CM$ 上诸复结构可积性的推论。在 \cite{Cecotti:1991me} 中，同样有一个模空间 $\vm$（它参数化 $d=2$ 的 $\CN=(2,2)$ 超对称场论），以及 $\vm$ 上的一个向量丛 $V$（即 Ramond 基态构成的丛）。人们在那里找到 $V$ 中一族由 $\zeta \in \IC^\times$ 参数化的平坦联络，其形式为
\begin{align}
\nabla_i &= \frac{1}{\zeta} C_i + D_i, \\
\nabla_{\bar i} &= \zeta \bar{C}_{\bar i} + \bar{D}_{\bar i},
\end{align}
其中 $D_i$ 是由耦合的绝热变化所提供的标准联络，诸 $C_i$ 则是「手征环」算符。这些联络的平坦性，是著名的 $tt^*$ 方程的推论。

在本文通篇，尤其是在第 5 节，许多构造——以及它们的物理解释——都与 $tt^*$ 故事中出现的那些相平行。% TODO 合稿改 \ref{sec:constructing-moduli}
